\subsection{问题一：静态游览路线规划}

问题一在全天开放时间内安排三类游客的路线。模型对象是项目访问序列，决策变量是项目集合、访问顺序和到达/开始/结束时刻。动态规划结果保存在 \verb|tab_q1_summary.csv| 和 \verb|tab_q1_routes.csv|。

\begin{figure}[H]
\centering
\includegraphics[width=0.88\linewidth]{fig_q1_model_schematic.png}
\caption{问题一动态规划建模流程。}
\end{figure}

\begin{figure}[H]
\centering
\includegraphics[width=0.9\linewidth]{fig_q1_result.png}
\caption{问题一不同游客和日期下的得分、排队和项目数。}
\end{figure}

\begin{figure}[H]
\centering
\includegraphics[width=0.9\linewidth]{fig_q1_validation.png}
\caption{问题一动态规划与贪心基线比较。}
\end{figure}

动态规划相对贪心基线平均得分提升 \(16.74\)。家庭亲子游工作日得分最高，为 \(96.734\)，完成 \(11\) 项非休整体验；节假日家庭亲子游在排队压力上升后仍完成 \(9\) 项非休整体验。
